\documentclass[10pt,a4paper]{amsart}
\usepackage[margin=19mm]{geometry}
\usepackage{amsmath,amssymb,mathtools}
\usepackage[scr=rsfs,cal=euler]{mathalfa}
\usepackage{xfrac}
\usepackage[unicode,hidelinks,bookmarks=false]{hyperref}
\newcommand{\ie}{i.e.\ }
\newcommand{\cf}{cf.\ }
\newcommand{\resp}{respectively\ }
\newcommand{\Subsection}{\subsection}
\providecommand{\nobreakdash}{\nobreak\hbox{-}}
\setlength{\emergencystretch}{2em}
\multlinegap0pt
\usepackage{amssymb}
\usepackage{tikz}
\newtheorem{lemma}{Lemma}
\title{Abstract Regular Polytopes of \\ Finite Irreducible Coxeter Groups}
\author{Malcolm Hoong Wai Chen, Peter Rowley}
\date{Source: arXiv:2501.01288v1 (CC BY 4.0)}
\begin{document}
\maketitle

\noindent\emph{Source and licence.} Abstract Regular Polytopes of \\ Finite Irreducible Coxeter Groups. Version arXiv:2501.01288v1 (CC BY 4.0), distributed by arXiv under the Creative Commons Attribution 4.0 International licence, http://creativecommons.org/licenses/by/4.0/, as recorded in the arXiv licence metadata for this exact version and shown on the abstract page https://arxiv.org/abs/2501.01288v1. Full licence notice: \texttt{licenses/arxiv-2501.01288-CC-BY-4.0.txt}.\par\smallskip

\noindent\textit{Editorial scope.} Counted unit: the article's lemma csr2 (source0.tex lines 455--457). The article's complete original proof of it is delivered with it (source0.tex lines 459--489, one \texttt{proof} environment of 31 lines). No mathematics is written, repaired or re-derived: the statement and the proof are the article's own lines, in the article's own notation, and no retained line is substituted or re-flowed.  Retained uncounted as the source-local context the proof needs: lines 299--301.  Nothing was omitted from any retained span, so the package carries no omission record.  No retained line contains a citation, so the wrapper carries no bibliography and no bibliography entry.  No retained line contains a figure environment, an image-inclusion command or drawing code, so no image is delivered and the package declares no asset dependency.  Editorial formatting only, in addition: an AMS \texttt{amsart} preamble that restates the article's own declarations and macros used by the retained lines, namely the environments \texttt{lemma} on separate counters, together with the standard packages those lines need.  The retained environments print in this document as Lemma 1 in order of appearance, whereas the article prints the same items with the numbering of its own full document.  The complete native source of the article as distributed in the arXiv e-print for this exact version is bundled unchanged as \texttt{sources/171058/source0.tex}.
\begin{lemma} \label{ipdn}
Let $G=\langle s_1,\dots, s_r\rangle \leq W \leq \mathrm{Sym}(2n)$ with $W \cong \mathrm{D}_n$ be a string group generated by involutions. If $G_{1,\dots,r-1}=S$ is a string C-group and $s_r \in N \setminus Z(G)$, then $G=W$ and is also a string C-group.
\end{lemma}

\begin{lemma} \label{csr2}
Let $n$ be even and $n \ge 6$. Then for every $4 \leq r \leq n-1$, $\{t_1,\dots,t_r\}$ is a rank $r$ C-string of $\mathrm{D}_n$ with Schl\"{a}fli type $\{3^{r-4},6,n-r+3,4\}$.
\end{lemma}

\begin{proof}
Each of the elements $t_1,\dots,t_r$ are involutions as they are defined as products of pairwise disjoint transpositions. Note that for every $1 \leq k \leq r-3$, we have
\begin{eqnarray*}
t_k t_r = 
\begin{cases}
(k,k+1)(n+k,n+k+1)(n-1,2n-1)(n,2n) \ &, \text{ if $r$ is odd } \\ 
(k,k+1)(n+k,n+k+1)(n-2,2n-2)(n-1,2n-1) \ &, \text{ if $r$ is even } \\
\end{cases},
\end{eqnarray*}
which are products of pairwise disjoint transpositions because 
\begin{eqnarray*}
    k+1 \leq (r-3)+1 \leq ((n-1)-3)+1 = n-3 < n-2.
\end{eqnarray*}
Likewise, if $r$ is odd we have
\begin{eqnarray*}
t_{r-2} t_r &=& (n-1,2n)(n,2n-1) (r-2,r-1)(r,r+1)\cdots(n-3,n-2) \\ &\phantom{=}& \quad(n+r-2,n+r-1)(n+r,n+r+1)\cdots(2n-3,2n-2),
\end{eqnarray*}
whereas if $r$ is even we have
\begin{eqnarray*}
t_{r-2} t_r &=& (n-2,2n-1)(n-1,2n-2)(r-2,r-1)(r,r+1)\cdots(n-4,n-3) \\ &\phantom{=}&\quad (n+r-2,n+r-1)(n+r,n+r+1)\cdots(2n-4,2n-3),
\end{eqnarray*}
which are also products of pairwise disjoint transpositions. Finally, we have
\begin{eqnarray*}
t_{r-1} t_r &=& (n-1,n-2,2n-1,2n-2)(n,2n) \\ &\phantom{=}& \quad(r-1,r)(r+1,r+2)\cdots(n-4,n-3) \\ &\phantom{=}& \quad (n+r-1,n+r)(n+r+1,n+r+2)\cdots(2n-4,2n-3)
\end{eqnarray*}
for $r$ odd and we also have
\begin{eqnarray*}
t_{r-1} t_r &=& (n-2,n-3,2n-2,2n-3)(n-1,n,2n-1,2n) \\ &\phantom{=}&\quad (r-1,r)(r+1,r+2)\cdots(n-5,n-4) \\ &\phantom{=}&\quad (n+r-1,n+r)(n+r+1,n+r+2)\cdots(2n-5,2n-4),
\end{eqnarray*}
for $r$ even. It follows from Lemma \ref{ipdn} that $\langle t_1,\dots,t_r \rangle=\mathrm{D}_n$ is a string C-group with Schl\"{a}fli type $\{3^{r-4},6,n-r+3,4\}$.
\end{proof}

\end{document}
